% Cover letter for the submission of paper IV (Ocelotas) to Research in Number Theory, 12 October 2026.
% Compile with: tectonic -X compile cover_letter.tex
\documentclass[11pt]{article}
\usepackage[margin=2.5cm]{geometry}
\usepackage{xurl}
\usepackage[hidelinks]{hyperref}
\setlength{\parindent}{0pt}
\setlength{\parskip}{0.8em}
\pagestyle{empty}

\begin{document}

\hfill 12 October 2026

To the Editors of \emph{Research in Number Theory}

Dear Editors,

I submit for consideration the manuscript ``On Terras's coefficient stopping time conjecture for Collatz sequences with few local maxima''.

Terras's coefficient stopping time (CST) conjecture says that a Collatz orbit descends below its starting value at the first step at which the coefficient $3^{q}/2^{k}$ drops below $1$. Rozier and Terracol (Discrete Math.\ 349, 2026, Remark~5.5) suggested studying it by the number of local maxima of the sequence. The paper does this by transferring the method of Simons and de Weger for $m$-cycles, based on the lower bound of Laurent, Mignotte and Nesterenko for linear forms in two logarithms, to the first crossing. It proves the conjecture for sequences with at most $24$ local maxima before the coefficient stopping time; with at most $81$ if the verification of Rozier and Terracol up to $2.8\cdot10^{19}$ is added; and with at most $42$ if instead our own computation up to $3\cdot10^{11}$ is used. The paper does not claim anything about the Collatz conjecture itself, and Section~1 states what is not proved.

Every step, including the evaluation of the finite certificates, is formalized in Lean~4 with Mathlib, with the external results as explicit hypotheses and only the standard axioms. The work was produced with extensive assistance from large language models, as documented in the section ``Use of AI tools''; that is why every step is either formalized or checked by a separate program.

For transparency:
\begin{itemize}
\item The manuscript is not under consideration at any journal. It has not been posted on a preprint server; an earlier draft is visible in the public repository mentioned below. On 12 October 2026 I sent an earlier draft privately to O.~Rozier, C.~Terracol and C.~Hercher, asking for their opinion; I have not received an answer yet.
\item A different paper of mine, on a conjecture of Gao on consecutive integers, is under consideration at the Journal of Automated Reasoning. It does not overlap with this one.
\item The sources, certificates and checking programs are in the public repository given in the data availability statement.
\end{itemize}

I am an independent researcher in Calama, Chile, without an academic affiliation.

Thank you for your consideration.

Sincerely,

Omar Javier Said Duran\\
Independent researcher, Calama, Chile\\
ORCID: \url{https://orcid.org/0009-0009-1418-2558}\\
\texttt{omarsaid24@gmail.com}

\end{document}
